\section{Experimental Setup}
\label{sec:experiments}

\subsection{Reproducible workload-drift benchmark}
To study continual index selection without depending on a particular DBMS build,
we use a \emph{cost-model} environment that is deterministic and fast while
retaining the structure of the real problem: a query stream organized into
workload phases, an index budget, per-column selectivity, and a reward defined as
the negative operating cost \eqref{eq:reward}. The environment implements the MDP
of Section~\ref{sec:problem} with the two-value latency model
\eqref{eq:latency}; in a production deployment, \(L(q,I)\) is replaced by a
measured latency or plan cost, and no other component changes.

\paragraph{Conflicting phases.}
The distinguishing feature of the benchmark is that phases \emph{conflict}: the
same query templates appear in several phases with \emph{opposite} optimal
actions, because a predicate that is highly selective in one phase is
non-selective in another. For example, \texttt{WHERE category=...} is worth
indexing in phase~1 and worth \emph{not} indexing in phase~2. The two phases are
therefore observationally near-identical except for the workload-context signal,
which makes the sequence strongly forgetting-prone---a database analogue of the
permuted-task benchmarks used in continual learning. Table~\ref{tab:phases}
lists the phases; a $\star$ marks a phase whose optimal actions conflict with its
partner.

\begin{table}[t]\centering\small
\caption{Workload phases. Phases $2,4,5$ share queries with phases $1,3,0$
respectively but invert the selectivity of the predicate columns, so the optimal
action for an identical observation flips.}
\label{tab:phases}
\begin{tabular}{cll}
\toprule
\# & Pattern & Selective columns \\
\midrule
0 & OLTP point lookups & \texttt{customer\_id}, \texttt{product\_id} \\
1 & Attribute filters & \texttt{category}, \texttt{status} \\
2$^\star$ & Filters (conflict) & none (inverse of \#1) \\
3 & Analytical scans & \texttt{order\_date}, \texttt{region} \\
4$^\star$ & Analytical (conflict) & none (inverse of \#3) \\
5$^\star$ & OLTP (conflict) & none (inverse of \#0) \\
6 & Return to \#0 & \texttt{customer\_id}, \texttt{product\_id} \\
\bottomrule
\end{tabular}
\end{table}

\subsection{Agents, baselines and ablations}
We compare three learning agents and three non-learning baselines.
\begin{itemize}
  \item \textbf{CLDB-RL} (ours): DQN $+$ drift-triggered EWC $+$
    database-aware prioritized replay $+$ action masking.
  \item \textbf{RL (no EWC)}: identical but $\lambda=0$; isolates the EWC effect.
  \item \textbf{RL (online, no replay)}: a replay buffer of capacity 64, i.e.\
    near-online fine-tuning; isolates the replay effect and serves as the
    catastrophic-forgetting control.
  \item \textbf{KEEP}, \textbf{RANDOM}, \textbf{GREEDY-CREATE}: static
    non-learning policies.
\end{itemize}

\subsection{Protocol}
For each seed we stream the seven phases in order. Within each phase the agent
acts \(\epsilon\)-greedily (masked), writing transitions to the buffer and
performing a DQN update every other step; the drift detector runs on the query
embeddings and triggers EWC consolidation when it fires. After each phase we
evaluate the greedy policy on every phase seen so far, filling one row of the
matrix \(\R\). Seeds are fixed 0--7 for reproducibility; we report the mean and
standard deviation over seeds and paired $t$-tests between the full model and
its ablations. All runs are single-threaded on a CPU.

\subsection{Hyperparameters}
Table~\ref{tab:hyper} lists the configuration. The embedding is produced offline
(the encoder is frozen); only the DQN is trained.

\begin{table}[t]\centering\small
\caption{Hyperparameters.}
\label{tab:hyper}
\begin{tabular}{ll}
\toprule
Component & Value \\
\midrule
State dim. & $74+384+3+3=464$ \\
Q-network & MLP $464\!\to\!128\!\to\!64\!\to\!3$, ReLU \\
Optimizer & Adam, lr $5\!\times\!10^{-4}$ \\
Discount $\gamma$ & $0.97$ \\
Batch size & $64$ \\
Target update & every $150$ grad.\ steps \\
$\epsilon$ schedule & $1.0\to0.05$ over $2500$ steps \\
Grad.\ clip & $10$ \\
Replay capacity & $20000$ (online variant: $64$) \\
Replay $\alpha,\beta$ & $0.6,\ 0.4$ \\
EWC $\lambda$ & $0.4$ \\
Index budget $B$ & $2$ \\
Build cost $\kappa$ & $0.5$ \\
Resource cost $\rho$ & $0.2$ \\
Latency $\ell_{\text{low}},\ell_{\text{high}}$ & $0.1,\ 1.0$ \\
Horizon (steps/phase) & $36$ \\
Episodes per phase & $12$ \\
Seeds & $8$ \\
\bottomrule
\end{tabular}
\end{table}
